\documentclass[11pt]{article}
\usepackage[margin=1in]{geometry}
\usepackage{booktabs}
\usepackage{hyperref}
\usepackage{xcolor}
\usepackage{enumitem}

\definecolor{reviewercolor}{RGB}{0,0,139}
\definecolor{responsecolor}{RGB}{0,100,0}

\newcommand{\reviewer}[1]{\textcolor{reviewercolor}{\textit{#1}}}
\newcommand{\response}[1]{\textcolor{responsecolor}{#1}}

\title{Response to Reviewers\\[0.5em]
\large ``When AI Replicates Science: A Meta-Study on Large Language Model Agents Reproducing Machine Learning Research''\\[0.5em]
\normalsize Revision 2.0}

\author{Claude Code and Fredrik Ahlgren}
\date{January 2025}

\begin{document}

\maketitle

\section*{Dear Reviewer,}

Thank you for your thorough and constructive review. Your critique was fair, substantive, and has significantly strengthened our work. We have undertaken a major revision addressing all concerns. Below, we provide point-by-point responses to your questions and detail the changes made.

\section*{Summary of Major Changes}

\begin{enumerate}
    \item \textbf{Terminology}: Reframed throughout from ``replication'' to ``documentation-driven re-enactment''
    \item \textbf{Uncertainty quantification}: All results now report mean $\pm$ std from 10 random seeds with bootstrap confidence intervals
    \item \textbf{Sensitivity analysis}: New analysis showing model rankings robust to generator parameter variations
    \item \textbf{Agent architecture}: Full specification including tools, decision log, and human intervention details
    \item \textbf{Related work}: New section engaging SciReplicate-Bench, LMR-BENCH, and xKG
    \item \textbf{AI Replicability Checklist}: Concrete documentation standards proposed
    \item \textbf{Self-replication framing}: Emphasized unique perspective as original author supervising AI re-enactment
\end{enumerate}

\section*{Detailed Responses to Reviewer Questions}

\subsection*{Question 1: Agent Architecture}

\reviewer{``What, precisely, was the agent architecture (tools, prompts, planning algorithm, allowed web/code access), and what interventions did the human supervisor perform? Can you provide a reproducible protocol and logs?''}

\response{We have added:
\begin{itemize}
    \item Section III.B specifying the agent system (Claude Code, claude-opus-4-5-20251101)
    \item Tools available: file read/write, shell execution, web search
    \item Human supervisor role: initial prompt, 2 checkpoint approvals, final review ($\sim$15 min total)
    \item Supplementary document \texttt{docs/agent\_architecture.md} with full decision log
    \item Supplementary document \texttt{docs/reproducibility\_protocol.md} with execution protocol
\end{itemize}
}

\subsection*{Question 2: Data Generator Release}

\reviewer{``Can you release the full synthetic data generator (code, parameters, random seeds) and provide sensitivity analyses demonstrating that your model rankings and R$^2$ claims are robust to reasonable variations in the generator?''}

\response{Yes. We have:
\begin{itemize}
    \item Added explicit equations in Section III.C (Equations 1-2)
    \item Added Table II validating synthetic data against documented statistics
    \item Created \texttt{experiments/sensitivity\_analysis.py} varying noise ($\pm$50\%), coefficients ($\pm$20\%), operating fractions ($\pm$10\%)
    \item Added Table IV showing model rankings remain stable across all variations
    \item Full generator code available with documented parameters and seeds
\end{itemize}
}

\subsection*{Question 3: Time-Series CV Procedure}

\reviewer{``Which model(s) correspond to the R$^2$ values reported in the time-series validation table, and what was the exact TSCV procedure (fold sizes, temporal ordering, nested tuning)?''}

\response{Section III now specifies:
\begin{itemize}
    \item Model: Random Forest (for comparison with random split)
    \item 5-fold expanding window using sklearn's TimeSeriesSplit
    \item Fold sizes: 5,000/10,000/15,000/20,000/25,000 train samples
    \item Test size: 5,000 samples per fold
    \item Limitation acknowledged: No nested CV for hyperparameter tuning
    \item Statistical significance test added (paired t-test, effect size)
\end{itemize}
}

\subsection*{Question 4: Sensitivity to Generator}

\reviewer{``How do your conclusions (e.g., `polynomial features outperform AutoML') change when the generator enforces different nonlinearities or noise structures that plausibly reflect ship dynamics? Please provide ablations.''}

\response{New Table IV reports sensitivity analysis:
\begin{itemize}
    \item Noise $\pm$50\%: R$^2$ varies by $\pm$0.003, rankings stable
    \item Coefficients $\pm$20\%: R$^2$ varies by $\pm$0.005, rankings stable
    \item Operating fraction $\pm$10\%: R$^2$ varies by $\pm$0.002, rankings stable
\end{itemize}
Ridge+Poly(d=2) remains top-ranked across all variations. Full results in supplementary materials.
}

\subsection*{Question 5: AutoGluon Comparison}

\reviewer{``Why not evaluate an up-to-date AutoML framework (e.g., AutoGluon) against TPOT to substantiate claims about `method evolution' and situate the 2018 approach in today's landscape?''}

\response{We attempted AutoGluon but encountered installation conflicts on the test environment. We have:
\begin{itemize}
    \item Noted AutoGluon as optional comparison in methodology
    \item Documented the installation issues in supplementary materials
    \item Focused comparison on methods that reliably executed across all seeds
    \item Acknowledged this as a limitation
\end{itemize}
The core finding---that polynomial features outperform complex ensembles---does not depend on AutoGluon comparison.
}

\subsection*{Question 6: Relation to SciReplicate-Bench}

\reviewer{``How do your findings relate to the negative results and difficulties reported in SciReplicate-Bench and LMR-BENCH for agentic reproduction? What problem class characteristics make your task feasible in `under one hour'?''}

\response{New Section II.D addresses this directly:
\begin{itemize}
    \item SciReplicate-Bench: Best LLM achieves 39\% execution accuracy
    \item LMR-BENCH: Reports ``persistent limitations in scientific reasoning''
    \item Our 100\% success requires explanation
\end{itemize}
We hypothesize task characteristics enabled success:
\begin{itemize}
    \item Small feature space (4-8 features vs. complex NLP)
    \item Interpretable physics (engine $\rightarrow$ fuel relationship)
    \item Available code repository with clear structure
    \item Well-documented statistics
\end{itemize}
This suggests AI re-enactment is not uniformly feasible---problem characteristics matter significantly.
}

\subsection*{Question 7: Taxonomy Classification}

\reviewer{``Can you explicitly classify your study using a standard taxonomy (repeatability, reproducibility, replicability; R1-R4-style categories) and clarify which dimension(s) your results support?''}

\response{New Section I.C defines our taxonomy (citing Barba 2018):
\begin{itemize}
    \item \textbf{Reproducibility}: Same data, same code $\rightarrow$ same results
    \item \textbf{Replicability}: New data, same methods $\rightarrow$ consistent results
    \item \textbf{Re-enactment}: Synthetic data from documentation $\rightarrow$ tests documentation sufficiency
\end{itemize}
We explicitly state: ``This study performs \textit{documentation-driven re-enactment}... a distinct but valuable question from whether original empirical claims are reproducible.''
}

\subsection*{Question 8: Contact Original Authors}

\reviewer{``Did you attempt to contact the original authors for missing details or to validate your inferences about preprocessing and data handling?''}

\response{This is a unique case: \textbf{the human supervisor (Fredrik Ahlgren) IS the original paper's first author}. We have emphasized this in Section I.D as a methodological strength:
\begin{itemize}
    \item Eliminates need to ``contact original authors''
    \item Provides insider perspective on documentation gaps
    \item Tests whether even original authors can reproduce from their own documentation
\end{itemize}
This self-replication perspective reveals documentation gaps invisible to authors but detectable through systematic AI analysis.
}

\subsection*{Question 9: Hardware/Software Environment}

\reviewer{``What hardware and software environment were used (CPU/GPU, library versions), and were multiple seeds evaluated to quantify variance?''}

\response{New supplementary document \texttt{docs/reproducibility\_protocol.md} specifies:
\begin{itemize}
    \item Platform: macOS (Darwin), ARM64 (Apple Silicon)
    \item Python 3.10+, scikit-learn $\geq$1.3.0, numpy $\geq$1.24.0
    \item GPU: Not used (CPU-only training)
    \item 10 random seeds evaluated: \{42, 123, 456, 789, 1024, 2048, 3141, 4269, 5555, 6789\}
    \item All results now report mean $\pm$ std with 95\% bootstrap CI
\end{itemize}
}

\subsection*{Question 10: AI Replicability Checklist}

\reviewer{``Do you envision a generalized, auditable checklist/tooling that operationalizes `AI replicability' of a paper's documentation? If so, what minimal items did this case study reveal?''}

\response{Yes. New Section V.C proposes a 6-item checklist based on gaps identified:
\begin{enumerate}
    \item \textbf{Seeds}: All random states explicitly documented
    \item \textbf{Preprocessing}: Data cleaning steps with code or pseudocode
    \item \textbf{Hyperparameters}: Sources stated (defaults, tuned, literature)
    \item \textbf{Statistics}: Sufficient for synthetic validation (mean, std, range, distribution shape)
    \item \textbf{Temporal structure}: Sampling interval, temporal dependencies noted
    \item \textbf{Negative handling}: Treatment of missing/invalid values specified
\end{enumerate}
We propose this as a concrete, actionable contribution toward ``AI replicability'' standards.
}

\section*{Response to Overall Assessment}

\reviewer{``The paper raises an important and timely question... However, the empirical foundation is too weak for the strength of the claims... As it stands, the work is better framed as a conceptual position piece with a suggestive, illustrative case study rather than as a definitive demonstration of AI-driven replication.''}

\response{We agree with this assessment and have revised accordingly:
\begin{itemize}
    \item Reframed as ``documentation-driven re-enactment'' (not ``replication'')
    \item Tempered claims about generalizability
    \item Explicitly acknowledged synthetic data limitations
    \item Positioned relative to SciReplicate-Bench difficulties
    \item Emphasized the contribution is methodological framework + checklist, not definitive proof
\end{itemize}
We believe the revised paper is now appropriately calibrated to its epistemic scope while still making a valuable contribution to the discourse on AI-driven research verification.
}

\section*{List of Changes by Section}

\begin{table}[h]
\centering
\begin{tabular}{ll}
\toprule
\textbf{Section} & \textbf{Changes} \\
\midrule
Abstract & Rewritten with uncertainty bounds, ``re-enactment'' framing \\
I.C (NEW) & Terminology taxonomy \\
I.D & Self-replication perspective emphasized \\
II.D (NEW) & AI-Driven Scientific Reproduction (3 new citations) \\
III.B (NEW) & Agent Architecture specification \\
III.C & Data generator equations (Eq. 1-2), validation table \\
III.D (NEW) & Uncertainty quantification methodology \\
Table II (NEW) & Synthetic data validation \\
Table III & Updated with mean $\pm$ std \\
Table IV (NEW) & Sensitivity analysis \\
V.C (NEW) & Benchmark comparison \\
V.D (NEW) & AI Replicability Checklist \\
VI & Conclusion rewritten with tempered claims \\
References & 4 new citations added \\
\bottomrule
\end{tabular}
\end{table}

\section*{Supplementary Materials Added}

\begin{itemize}
    \item \texttt{docs/agent\_architecture.md} - Full agent specification
    \item \texttt{docs/reproducibility\_protocol.md} - Environment and execution protocol
    \item \texttt{experiments/run\_uncertainty\_analysis.py} - Multi-seed experiments
    \item \texttt{experiments/sensitivity\_analysis.py} - Generator robustness tests
\end{itemize}

\vspace{1em}

We thank the reviewer again for their thorough and constructive feedback. We believe these revisions substantially strengthen the paper and address all major concerns.

\vspace{2em}

\noindent Sincerely,\\[1em]
Claude Code and Fredrik Ahlgren

\end{document}
